\begin{mistake}{2026-05-25}{03}

\begin{problemcontent}
设 \( f(x) \) 在 \( (-\infty, +\infty) \) 有定义，\( f(0) = 1 \) 且 \( f(x) \) 在 \( x = 0 \) 处连续，并满足 \( f(2x) = f(x) \)，则在 \( (-\infty, +\infty) \) 上 \( f(x) = \) \underline{\quad 1 \quad}。
\end{problemcontent}

\begin{solutioncontent}
根据已知条件 \( f(2x) = f(x) \)，等价变形为 \( f(x) = f\left(\frac{x}{2}\right) \)。
对于任意实数 \( x \)，不断递推可得：
\[ f(x) = f\left(\frac{x}{2}\right) = f\left(\frac{x}{2^2}\right) = \cdots = f\left(\frac{x}{2^n}\right) \]
对等式两边令 \( n \to \infty \) 取极限：
\[ f(x) = \lim_{n \to \infty} f\left(\frac{x}{2^n}\right) \]
因为对于任意实数 \( x \)，均有 \( \lim_{n\to\infty} \frac{x}{2^n} = 0 \)，且已知 \( f(x) \) 在 \( x=0 \) 处连续，根据连续函数的性质，极限符号可移入函数内部：
\[ \lim_{n \to \infty} f\left(\frac{x}{2^n}\right) = f\left( \lim_{n \to \infty} \frac{x}{2^n} \right) = f(0) \]
代入已知条件 \( f(0) = 1 \)，即得：
\[ f(x) = 1 \]
由于 \( x \) 是任意选取，故在 \( (-\infty, +\infty) \) 上 \( f(x) \equiv 1 \)。
\end{solutioncontent}

\mistakeknowledge{
\begin{itemize}
  \item \textbf{核心方法}：处理 \( f(kx)=f(x) \) 形式的方程，通过递推将自变量缩放到已知连续的点（如 \( x=0 \)）。
  \item \textbf{连续性应用}：若 \( f(x) \) 在 \( x_0 \) 连续且 \( \lim x_n = x_0 \)，则 \( \lim f(x_n) = f(\lim x_n) = f(x_0) \)。
  \item \textbf{易错点}：必须依赖连续性才能将极限符号移入函数内，否则无法求出具体值。
\end{itemize}}

\end{mistake}
